\section{Every logarithmic coefficient correction at order three}
\label{sec:logarithmic}

In this section $s=3$, $u=\Uc(3)$, $A=\zeta(2)$, and
$a_n=a_n(3,\Uc(3))$.  The exact Lambert core from
\cref{thm:chart-three} determines all inverse powers of the large
logarithm in the leading singular sector.  The analytic feedback
correction satisfies $x-\tau=O(\tau^2)$ and is smaller than every
such logarithmic correction relative to $\tau$.

\subsection{A convergent logarithmic expansion of the exact core}

Put
\begin{equation}\label{eq:log-core-parameters}
 L=\log(1/d),\qquad c=5-\log(4A),\qquad T=2\Lambda.
\end{equation}
The identities defining $\Lambda$ are equivalent to
\begin{equation}\label{eq:log-T}
 T-\log T=L+c,\qquad \tau=2\sqrt{Ad/T}.
\end{equation}
Both the factor $4A$ and the additive constant $5$ enter the final
coefficient correction.

\begin{proposition}[All core logarithmic polynomials]
\label{prop:core-log-polynomials}
For sufficiently small positive $d$,
\begin{equation}\label{eq:core-log-series}
 \tau=2\sqrt{Ad}\,L^{-1/2}
 \sum_{j\ge0}\frac{P_j(\log L+c)}{L^j},
\end{equation}
where this series converges, and
\begin{equation}\label{eq:core-log-polynomial-formula}
 P_0(Y)=1,\qquad
 P_j(Y)=-\frac1{2j}[h^{j-1}]
 (1+h)^{-3/2}\{Y+\log(1+h)\}^{j}\quad(j\ge1).
\end{equation}
Each $P_j$ has degree $j$.  For every fixed $M\ge0$, truncation
after $j=M$ has relative remainder
$O((\log L)^{M+1}/L^{M+1})$.
\end{proposition}

\begin{proof}
Write $T=L(1+V)$.  Equation \eqref{eq:log-T} becomes
\[
 V=\eta+p\log(1+V),\qquad
 p=L^{-1},\quad \eta=(\log L+c)/L.
\]
The implicit derivative in $V$ at $(p,\eta,V)=(0,0,0)$ is one.
Thus $V$ and $(1+V)^{-1/2}$ are holomorphic in the two small
coordinates $(p,\eta)$.  After setting $\eta=pY$, apply ordinary
Lagrange inversion to $V=p\{Y+\log(1+V)\}$ and the function
$(1+V)^{-1/2}$.  It gives
\eqref{eq:core-log-polynomial-formula}.
Grouping the normally convergent double series by total degree gives
\eqref{eq:core-log-series}; both $p$ and $pY$ are small on the physical
curve.  The leading coefficient of $P_j$ is
$(-1)^j\binom{2j}{j}/4^j$, which is nonzero.
Cauchy bounds for the double series give the stated remainder and its
uniform counterpart on closed subsectors.
\end{proof}

The first polynomials are
\begin{align}
 P_1(Y)&=-\frac Y2,\label{eq:core-log-P1}\\
 P_2(Y)&=\frac{3Y^2}{8}-\frac Y2,\label{eq:core-log-P2}\\
 P_3(Y)&=-\frac{5Y^3}{16}+Y^2-\frac Y2.\label{eq:core-log-P3}
\end{align}
In particular,
\begin{equation}\label{eq:three-x-logarithmic}
 x=2\sqrt{Ad/L}\left[
 1-\frac{\log L-\log(4A)+5}{2L}
 +O\left(\frac{(\log L)^2}{L^2}\right)\right].
\end{equation}
The $O(\tau^2)$ correction from \eqref{eq:three-chart-correction}
is absorbed here because it is smaller than every fixed inverse power
of $L$ relative to the leading term.

\subsection{A finite formula after singularity transfer}

It is useful to state the transfer calculation explicitly: otherwise
the shift from the Gamma factor can change the constant in the answer.
For $a$ near $1/2$ define
\begin{equation}\label{eq:transfer-multipliers}
 G(a)=\frac1{\Gamma(-a)},\qquad
 M_k(a)=\frac{(-1)^kG^{(k)}(a)}{k!G(a)}.
\end{equation}
For the finite sums of functions $L^{-\beta}P(\log L)$ used here,
power--logarithm transfer \cite{FO} can be expressed as the asymptotic
operator
\begin{equation}\label{eq:log-transfer-operator}
 [z^n](1-z)^aF\left(\log\frac1{1-z}\right)
 \sim\frac{n^{-a-1}}{\Gamma(-a)}
 \left.\frac{\Gamma(-a)}{\Gamma(-a+\partial_L)}F(L)
 \right|_{L=\log n},
\end{equation}
where the operator means
$\sum_{k\ge0}M_k(a)\partial_L^k$ as an asymptotic derivative
expansion.  No convergence of this infinite operator is claimed.
For a finite target order only finitely many derivatives contribute.

The signs can be checked directly on the Hankel contour used in
singularity analysis.  The substitution $1-z=t/n$ replaces the
logarithmic argument by $\log n-\log t$.  Taylor expansion of $F$
there contributes $(-1)^kF^{(k)}(\log n)(\log t)^k/k!$.
Its Hankel integral is $G^{(k)}(a)$, giving
\eqref{eq:log-transfer-operator}.  The analytic continuation from
\cref{lem:delta} permits the contour deformation.  Truncating the
logarithmic expansion and its differentiated versions gives the usual
uniform remainder estimates; powers of $1/n$ from the contour scaling
are smaller than every fixed inverse power of $\log n$.

\begin{theorem}[All logarithmic coefficient corrections]
\label{thm:all-logarithmic-coefficients}
For $N\ge0$ define the polynomial
\begin{equation}\label{eq:transferred-log-polynomials}
 \mathcal Q_N(Y)=\sum_{j=0}^{N}M_{N-j}(1/2)
 \prod_{\ell=0}^{N-j-1}
 \left(\frac{\dd}{\dd Y}-j-\frac12-\ell\right)P_j(Y),
\end{equation}
where the empty product is one.  For every fixed $M\ge0$,
\begin{equation}\label{eq:all-log-coefficient-law}
 \begin{split}
 \Rb^n a_n={}&\sqrt{\frac A\pi}\,
 n^{-3/2}(\log n)^{-1/2}\left[
 \sum_{N=0}^{M}
 \frac{\mathcal Q_N(\log\log n+c)}{(\log n)^N}
 \right.\\[-2pt]
 &\left.\hspace{35mm}
 +O\left(\frac{(\log\log n)^{M+1}}{(\log n)^{M+1}}\right)
 \right].
 \end{split}
\end{equation}
Thus every correction is supplied by finite polynomial and
Gamma-derivative operations.
\end{theorem}

\begin{proof}
For the singular part of $v=\vc+d-x$, replace $x$ by the exact
core $\tau$.  The omitted terms $x-\tau=O(\tau^2)$ contribute at
coefficient orders smaller by a positive power of $n$, up to powers of
its logarithm.  Locally analytic terms do not affect the singular
expansion.  Also $d=-\Log(r/\Rb)=1-r/\Rb+O((1-r/\Rb)^2)$;
this replacement changes only corrections smaller by a power of $n$.

Apply \eqref{eq:log-transfer-operator}, with $a=1/2$, to each
term $L^{-j-1/2}P_j(\log L+c)$ in
\eqref{eq:core-log-series}.  Repeated differentiation gives exactly
\[
 \partial_L^k\{L^{-j-1/2}P_j(\log L+c)\}
 =L^{-j-k-1/2}
 \prod_{\ell=0}^{k-1}
 \left(\frac{\dd}{\dd Y}-j-\frac12-\ell\right)P_j(Y)
 \bigg|_{Y=\log L+c}.
\]
Collecting $j+k=N$ proves \eqref{eq:transferred-log-polynomials}.
The amplitude of $-\tau$ is $-2\sqrt A$, and
$\Gamma(-1/2)=-2\sqrt\pi$, giving the positive prefactor in
\eqref{eq:all-log-coefficient-law}.  The finite-order transfer
remainder and \cref{prop:core-log-polynomials} prove the stated error.
\end{proof}

\subsection{The first three terms in closed form}

Let $\psi=\Gamma'/\Gamma$, $\psi_1=\psi'$, and abbreviate their
values at $-1/2$ by $\psi_*$ and $\psi_{1,*}$.  Formula
\eqref{eq:transfer-multipliers} gives
\[
 M_0=1,\qquad M_1=-\psi_*,\qquad
 M_2=\frac{\psi_*^2-\psi_{1,*}}2.
\]
Substituting \eqref{eq:core-log-P1}--\eqref{eq:core-log-P2} into
\eqref{eq:transferred-log-polynomials} yields
\begin{align*}
 \mathcal Q_0(Y)&=1,\\
 \mathcal Q_1(Y)&=-\frac{Y-\psi_*}{2},\\
 \mathcal Q_2(Y)&=\frac38(Y-\psi_*)^2-\frac12(Y-\psi_*)
                  -\frac38\psi_{1,*}.
\end{align*}
The standard digamma recurrence and half-integer values give
\[
 \psi_*=2-\gamma-2\log2,\qquad
 \psi_{1,*}=\frac{\pi^2}{2}+4.
\]
Consequently, with
\begin{equation}\label{eq:log-Hn}
 H_n=\log\log n-\log A+\gamma+3,
\end{equation}
we obtain the explicit refinement
\begin{equation}\label{eq:three-term-log-law}
 \boxed{\begin{aligned}
 \Rb^n a_n={}&\sqrt{\frac A\pi}\,
 n^{-3/2}(\log n)^{-1/2}\left[
 1-\frac{H_n}{2\log n}\right.\\
 &\left.\quad+
 \frac{\frac38H_n^2-\frac12H_n-\frac{3\pi^2}{16}-\frac32}
 {(\log n)^2}
 +O\left(\frac{(\log\log n)^3}{(\log n)^3}\right)
 \right].
 \end{aligned}}
\end{equation}
The cancellation of $\log4$ between the exact core and the digamma
shift leaves the simple constant $-\log A+\gamma+3$.
The core and feedback charts are convergent analytic constructions;
\eqref{eq:all-log-coefficient-law} and
\eqref{eq:three-term-log-law} are coefficient asymptotic expansions.
These are distinct convergence statements.
